% GENERATED FILE. Edit canonical lessons, then run npm run generate:books.
% canonical-source-hash: fnv1a64:c9733a1b7e974364
% canonical-lessons: TE-C392-mulam, TE-C392-udbhavincu, TE-C392-daritiyu, TE-C392-kalugu, TE-C392-talettu

\chapter{Where Things Come From}
\label{ch:te-where-things-come-from}
\hlchaptermodality{392}

\begin{chapteropening}
\textbf{By the end of this chapter:} \emph{I can talk about sources, origins and what something leads to.}

Complete the last lesson of chapter 392: I can talk about sources, origins and what something leads to.
\end{chapteropening}

\section[\texorpdfstring{\te{మూలం} (mūlaṁ)}{మూలం (mūlaṁ)}]{\te{మూలం} (mūlaṁ) — a source; an origin}

\label{lesson:TE-C392-mulam}



\begin{quote}
Before the new one: say the Telugu for to stress, then the Telugu for to bring up.
\end{quote}

\subsection*{You'll want to know: \te{మూలం}}
\textbf{\te{మూలం}} — \emph{mūlaṁ} — \textquotedblleft{}a source; an origin\textquotedblright{}.

\begin{quote}
\te{ఆ} \te{వార్తకు} \te{మూలం} \te{ఏమిటి}?
\end{quote}

\begin{quote}
\emph{ā vārtaku mūlaṁ ēmiṭi?}
\end{quote}

\begin{quote}
What is the source of that news?
\end{quote}

First a root: \te{మూల} \te{కారణం} is the root cause.

\begin{cousinweb}
\textbf{\te{మూలం}} (\emph{mūlaṁ}) is borrowed from Sanskrit \emph{mūla}, \textquotedblleft{}a root\textquotedblright{} (CDIAL 10250); in Telugu it is the root of anything, its source or origin.
\end{cousinweb}

\begin{grammarlens}[title={What you've built}]
One more word for giving reasons and opinions.
\end{grammarlens}

\subsection*{Your turn}
\begin{itemize}
\raggedright
  \item \emph{Say it:} \emph{mūlaṁ}
  \item \emph{Say it:} \emph{mūlaṁ}, once more
  \item \emph{Say it:} \emph{ā vārtaku mūlaṁ ēmiṭi?} — What is the source of that news?
  \item \emph{Recall:} say the Telugu for to stress, then the Telugu for to bring up, then say \emph{mūlaṁ} again
\end{itemize}

\begin{tcolorbox}[breakable,colback=teal!4,colframe=teal!35!black,title={Before you move on}]
What does \te{మూలం} mean? (\textquotedblleft{}A source; an origin\textquotedblright{}.) Say it once more.
\end{tcolorbox}
\section[\texorpdfstring{\te{ఉద్భవించు} (udbhaviñcu)}{ఉద్భవించు (udbhaviñcu)}]{\te{ఉద్భవించు} (udbhaviñcu) — to originate; to spring from}

\label{lesson:TE-C392-udbhavincu}



\begin{quote}
Before the new one: say the Telugu for to bring up, then the Telugu for a source.
\end{quote}

\subsection*{You'll want to know: \te{ఉద్భవించు}}
\textbf{\te{ఉద్భవించు}} — \emph{udbhaviñcu} — \textquotedblleft{}to originate; to spring from\textquotedblright{}.

\begin{quote}
\te{ఆ} \te{ఆచారం} \te{చాలా} \te{పాత} \te{కాలంలో} \te{ఉద్భవించింది}.
\end{quote}

\begin{quote}
\emph{ā ācāraṁ cālā pāta kālaṁlō udbhaviñcindi.}
\end{quote}

\begin{quote}
That custom arose in very old times.
\end{quote}

Past (it …): \te{ఉద్భవించింది}, \emph{udbhaviñcindi}.

Formal: for customs, ideas, languages and rivers, so its past is shown as \te{ఉద్భవించింది}, it arose.

\begin{cousinweb}
\textbf{\te{ఉద్భవించు}} (\emph{udbhaviñcu}) is built on Sanskrit \emph{udbhava}, \textquotedblleft{}origin, birth\textquotedblright{} (CDIAL 2039), from the root \emph{bhū}, \textquotedblleft{}to become\textquotedblright{}, with the Telugu verb ending -\te{ఇంచు}.
\end{cousinweb}

\begin{grammarlens}[title={What you've built}]
One more word for giving reasons and opinions.
\end{grammarlens}

\subsection*{Your turn}
\begin{itemize}
\raggedright
  \item \emph{Say it:} \emph{udbhaviñcu}
  \item \emph{Say it:} \emph{udbhaviñcu}, once more
  \item \emph{Say it:} \emph{ā ācāraṁ cālā pāta kālaṁlō udbhaviñcindi.} — That custom arose in very old times.
  \item \emph{Recall:} say the Telugu for to bring up, then the Telugu for a source, then say \emph{udbhaviñcu} again
\end{itemize}

\begin{tcolorbox}[breakable,colback=teal!4,colframe=teal!35!black,title={Before you move on}]
What does \te{ఉద్భవించు} mean? (\textquotedblleft{}To originate; to spring from\textquotedblright{}.) Say it once more.
\end{tcolorbox}
\section[\texorpdfstring{\te{దారితీయు} (dāritīyu)}{దారితీయు (dāritīyu)}]{\te{దారితీయు} (dāritīyu) — to lead to; to result in}

\label{lesson:TE-C392-daritiyu}



\begin{quote}
Before the new one: say the Telugu for a source, then the Telugu for to originate.
\end{quote}

\subsection*{You'll want to know: \te{దారితీయు}}
\textbf{\te{దారితీయు}} — \emph{dāritīyu} — \textquotedblleft{}to lead to; to result in\textquotedblright{}.

\begin{quote}
\te{చిన్న} \te{అభిప్రాయభేదం} \te{పెద్ద} \te{గొడవకు} \te{దారితీసింది}.
\end{quote}

\begin{quote}
\emph{cinna abhiprāyabhēdaṁ pedda goḍavaku dāritīsindi.}
\end{quote}

\begin{quote}
A small difference of opinion led to a big quarrel.
\end{quote}

Past (it …): \te{దారితీసింది}, \emph{dāritīsindi}.

What it leads to takes -\te{కి}: \te{గొడవకు} \te{దారితీయు}. Word for word, to take the road. Events lead to things, so its past is shown as \te{దారితీసింది}, it led to.

\begin{cousinweb}
\textbf{\te{దారితీయు}} (\emph{dāritīyu}) is native Dravidian, a compound of \emph{dāri}, \textquotedblleft{}a way, a road\textquotedblright{} (DEDR 3170), and \emph{tīyu}, \textquotedblleft{}to take, to draw\textquotedblright{} (DEDR 3407).
\end{cousinweb}

\begin{grammarlens}[title={What you've built}]
One more word for giving reasons and opinions.
\end{grammarlens}

\subsection*{Your turn}
\begin{itemize}
\raggedright
  \item \emph{Say it:} \emph{dāritīyu}
  \item \emph{Say it:} \emph{dāritīyu}, once more
  \item \emph{Say it:} \emph{cinna abhiprāyabhēdaṁ pedda goḍavaku dāritīsindi.} — A small difference of opinion led to a big quarrel.
  \item \emph{Recall:} say the Telugu for a source, then the Telugu for to originate, then say \emph{dāritīyu} again
\end{itemize}

\begin{tcolorbox}[breakable,colback=teal!4,colframe=teal!35!black,title={Before you move on}]
What does \te{దారితీయు} mean? (\textquotedblleft{}To lead to; to result in\textquotedblright{}.) Say it once more.
\end{tcolorbox}
\section[\texorpdfstring{\te{కలుగు} (kalugu)}{కలుగు (kalugu)}]{\te{కలుగు} (kalugu) — to arise; to come about}

\label{lesson:TE-C392-kalugu}



\begin{quote}
Before the new one: say the Telugu for to originate, then the Telugu for to lead to.
\end{quote}

\subsection*{You'll want to know: \te{కలుగు}}
\textbf{\te{కలుగు}} — \emph{kalugu} — \textquotedblleft{}to arise; to come about\textquotedblright{}.

\begin{quote}
\te{ఆ} \te{వార్త} \te{విని} \te{నాకు} \te{ఆనందం} \te{కలిగింది}.
\end{quote}

\begin{quote}
\emph{ā vārta vini nāku ānandaṁ kaligindi.}
\end{quote}

\begin{quote}
Hearing that news, I felt joy.
\end{quote}

Past (it …): \te{కలిగింది}, \emph{kaligindi}.

The feeling is the subject and the person takes -\te{కి}: \te{నాకు} \te{సందేహం} \te{కలిగింది}, a doubt arose in me. So its past is shown in the third person, \te{కలిగింది}; no one says \te{కలిగాను}.

\begin{cousinweb}
\textbf{\te{కలుగు}} (\emph{kalugu}) is native Dravidian: DEDR 1300 lists Telugu \emph{kalugu}, \textquotedblleft{}to happen, to occur, to be produced, to exist\textquotedblright{}, with Tamil \emph{kala}, \textquotedblleft{}to appear, to come into being\textquotedblright{}.
\end{cousinweb}

\begin{grammarlens}[title={What you've built}]
One more word for giving reasons and opinions.
\end{grammarlens}

\subsection*{Your turn}
\begin{itemize}
\raggedright
  \item \emph{Say it:} \emph{kalugu}
  \item \emph{Say it:} \emph{kalugu}, once more
  \item \emph{Say it:} \emph{ā vārta vini nāku ānandaṁ kaligindi.} — Hearing that news, I felt joy.
  \item \emph{Recall:} say the Telugu for to originate, then the Telugu for to lead to, then say \emph{kalugu} again
\end{itemize}

\begin{tcolorbox}[breakable,colback=teal!4,colframe=teal!35!black,title={Before you move on}]
What does \te{కలుగు} mean? (\textquotedblleft{}To arise; to come about\textquotedblright{}.) Say it once more.
\end{tcolorbox}
\section[\texorpdfstring{\te{తలెత్తు} (talettu)}{తలెత్తు (talettu)}]{\te{తలెత్తు} (talettu) — to crop up; to surface}

\label{lesson:TE-C392-talettu}



\begin{quote}
Before the new one: say the Telugu for to lead to, then the Telugu for to arise.
\end{quote}

\subsection*{You'll want to know: \te{తలెత్తు}}
\textbf{\te{తలెత్తు}} — \emph{talettu} — \textquotedblleft{}to crop up; to surface\textquotedblright{}.

\begin{quote}
\te{కొత్త} \te{ఇంట్లో} \te{చాలా} \te{సమస్యలు} \te{తలెత్తాయి}.
\end{quote}

\begin{quote}
\emph{kotta iṇṭlō cālā samasyalu talettāyi.}
\end{quote}

\begin{quote}
A lot of problems cropped up in the new house.
\end{quote}

Past (it …): \te{తలెత్తింది}, \emph{talettindi}.

\te{తల}, the head, + \te{ఎత్తు}, to raise: a problem raises its head. In this sense things crop up, not people, so its past is shown as \te{తలెత్తింది}, it cropped up.

\begin{cousinweb}
\textbf{\te{తలెత్తు}} (\emph{talettu}) is native Dravidian, a compound of \emph{tala}, \textquotedblleft{}the head\textquotedblright{} (DEDR 3103), and \emph{ettu}, \textquotedblleft{}to lift, to raise\textquotedblright{} (DEDR 796).
\end{cousinweb}

\begin{grammarlens}[title={What you've built}]
One more word for giving reasons and opinions.
\end{grammarlens}

\subsection*{Your turn}
\begin{itemize}
\raggedright
  \item \emph{Say it:} \emph{talettu}
  \item \emph{Say it:} \emph{talettu}, once more
  \item \emph{Say it:} \emph{kotta iṇṭlō cālā samasyalu talettāyi.} — A lot of problems cropped up in the new house.
  \item \emph{Recall:} say the Telugu for to lead to, then the Telugu for to arise, then say \emph{talettu} again
\end{itemize}

\begin{tcolorbox}[breakable,colback=teal!4,colframe=teal!35!black,title={Before you move on}]
What does \te{తలెత్తు} mean? (\textquotedblleft{}To crop up; to surface\textquotedblright{}.) Say it once more.
\end{tcolorbox}
